% ===== Preâmbulo =====
% ----- Definição da classe de documento -----
\documentclass[twoside]{article}

% ----- Pacotes -----
\usepackage[utf8]{inputenc}
\usepackage[brazil]{babel}
\usepackage{geometry}
\usepackage{setspace}
\usepackage[font=small,labelfont=bf]{caption}
\usepackage{graphicx}
\usepackage{amsmath, amssymb, mathtools}
\usepackage{amssymb}
\usepackage[backend=biber]{biblatex}
\usepackage{enumitem}
\usepackage{pifont}
\usepackage{tabularx}
\usepackage{booktabs}
\usepackage{multirow}
\usepackage{caption}
\usepackage{subcaption}
\usepackage{multicol}
\usepackage{listings}
\usepackage[dvipsnames]{xcolor}
\usepackage{colortbl}
\usepackage{fancyhdr}
\usepackage{background}
\usepackage{pgfplots}
\usepackage{mhchem}
\usepackage{tikz}
\usepackage{circuitikz}
\usepackage{qrcode}
\usepackage[
colorlinks=true,
linkcolor=blue,
citecolor=red,
urlcolor=magenta,
pdfstartview=FitH,
bookmarks=true,
hidelinks
]{hyperref}

\usetikzlibrary{shapes.geometric, arrows.meta, positioning}

\usepackage{lipsum}

% ----- Configurações -----
\DeclareMathOperator{\sen}{sen}

\geometry{a4paper, margin=2.5cm}

\lstset{
	basicstyle=\ttfamily\small,
	keywordstyle=\color{blue},
	stringstyle=\color{orange},
	commentstyle=\color{gray},
	numbers=left,
	numberstyle=\tiny,
	stepnumber=1,
	numbersep=5pt,
	backgroundcolor=\color{gray!10},
	frame=single,
	captionpos=b
}

\pagestyle{fancy}
\fancyhead[L]{Exemplos de \LaTeX{}}
\fancyhead[C]{}
\fancyhead[R]{\today}
\fancyfoot[C]{\LaTeX{}}
\fancyfoot[lo]{\thepage}
\fancyfoot[ro]{UFPR}
\fancyfoot[le]{UFPR}
\fancyfoot[re]{\thepage}

\backgroundsetup{
	scale=8,
	color=gray,
	opacity=0.2,
	angle=45,
	position=current page.center,
	contents={}
}

% ----- Bibliografia -----
\bibliography{referencias.bib}

% ----- Definições de Título -----
\title{Exemplos de \LaTeX{}}
\author{Franco, Davi}
\date{\today}

% ===== Corpo =====
\begin{document}
	
% ----- Página de título -----
	\begin{titlepage}
		\backgroundsetup{contents={\LaTeX{} UFPR}}
		\BgThispage
		\begin{center}
			\vspace*{2cm}
			
			{\Large \textbf{UNIVERSIDADE FEDERAL DO PARANÁ}}\\[0.5cm]
			{\large Bacharelado em Física}\\[3cm]
			
			{\Huge \textbf{Exemplos de \LaTeX{}}}\\[2cm]
			
			{\Large Franco, Davi}\\[0.5cm]
			\vfill
			
			{\large Curitiba – PR\\
			\today}
		\end{center}
	\end{titlepage}
	
% ----- Contra-capa -----

	\maketitle
	\newpage
	
% ----- Sumário -----
	\tableofcontents
	\newpage
	
% ----- Ênfase e alinhamento -----
	\section{Ênfase e alinhamento}
	
	\subsection{Ênfase e Alinhado à direita}
	
	\begin{flushright}
		\emph{Ênfase, \lipsum[1]}
	\end{flushright}
	
	\subsection{Itálico e Alinhado à esquerda}
	
	\begin{flushleft}
		\textit{Itálico, \lipsum[2]}
	\end{flushleft}

	\subsection{Negrito e Centralizado}
	
	\begin{center}
		\textbf{Negrito, \lipsum[3]}
	\end{center}


	\subsection{Sublinhado}

	\underline{Sublinhado, este não obedece a margem.}

	\subsection{Tamanho de fonte}
	
	\begin{center}
		{\tiny Muito pequena} 

		{\small Pequena} 

		{\normalsize Normal} 

		{\large Grande} 

		{\huge Muito grande}

		{\fontsize{25}{10}\selectfont Tamanho personalizado}
	\end{center}
	
	\hrule
	
% ----- Seções e Subseções -----
	\section{Seções e Subseções}
	\lipsum[1]
	
	\subsection{subseção 1}
	\lipsum[2]
	
	\subsection{subseção 2}
	\lipsum[3-4]
	
	\subsection{subseção 3}
	\lipsum[5]
	
	\subsubsection{subsubseção}
	\lipsum[6]
	
% ----- Iteinização e Enumeração
	\section{Iteinização e Enumeração}
	
	\subsection{Inteinização}
	
	\begin{itemize}
		\item Primeiro item
		\item Segundo item
		\item Terceiro item
		\begin{itemize}
			\item Subitem A
			\begin{itemize}
				\item Subitem B
				\begin{itemize}
					\item Subitem C
				\end{itemize}
			\end{itemize}
		\end{itemize}
	\end{itemize}
	
	\subsubsection{Marcador personalizado}
	
	\begin{itemize}[label=→]
		\item Item com seta
	\end{itemize}
	
	Ou
	
	\begin{itemize}[label=\ding{72}]
		\item Item com estrela
	\end{itemize}

	\subsubsection{Marcadores diferentes}
	
		\begin{itemize}
			\item[] Primeiro item
			\item[La] Segundo item
			\item[TeX´] Terceiro item
		\end{itemize}
	
	\subsection{Enumeração}
	
	\begin{enumerate}
	\item Primeiro item
	\item Segundo item
	\item Terceiro item
		\begin{enumerate}
			\item Subitem A
			\begin{enumerate}
				\item Subitem B
				\begin{enumerate}
					\item Subitem C
				\end{enumerate}
			\end{enumerate}
		\end{enumerate}
	\end{enumerate}

	\subsection{Marcador personalizado}
	
	\begin{enumerate}[label=\alph*)]
		\item Letra a)
		\item Letra b)
	\end{enumerate}

	\begin{enumerate}[label=\roman*.]
		\item um.
		\item dois.
	\end{enumerate}

	\subsection{Começar a qualquer momento}
	
	\begin{enumerate}[start=5]
		\item Começa do item 5
	\end{enumerate}
		
	\subsection{Espaçamento personalizado}
	
	\begin{itemize}[leftmargin=2cm, itemsep=0.5cm]
		\item Primeiro item
		\item Segundo item
		\item Terceiro item
	\end{itemize}

% ----- Multiplas colunas -----
	\section{Múltiplas colunas}
	
	\begin{multicols}{2}
		\lipsum[1-2]
	\end{multicols}

	\begin{multicols}{3}
		\lipsum[3-4]
	\end{multicols}

% ----- Matemática -----
	\section{Matemática}
	
	\subsection{Equação inline}
	Curabitur dictum gravida mauris. Nam arcu libero, nonummy eget, consectetuer id, vulputate a, magna. Donec vehicula augue eu neque $\vec{F_{r}} = m \cdot \vec{a}$ p´ellentesque habitant morbi tristique senectus et netus et malesuada fames ac turpis egestas. Mauris ut leo $\vec{v} = \vec{v_0} + \vec{a} \cdot t$.
	
	\subsection{Equações em bloco}
	
	\subsubsection{Equação numerada}
	
	\begin{equation}
		F(x) = \int f(x) dx
		\label{eq:antideriv}
	\end{equation}
	
	\subsubsection{Equação não numerada}
	
	\begin{equation*}
		\pi = \int_{-1}^{1} \frac{dx}{\sqrt{1-x^2}}
	\end{equation*}

	\subsubsection{Equações alinhadas numeradas}

	\begin{align}
		\pi &= \int_{-1}^{1} \frac{dx}{\sqrt{1-x^2}}	\label{eq:pi} \\
		F(x) &= \int f(x) dx
		\label{eq:int}
	\end{align}

	\subsubsection{Equações alinhadas não numeradas}
	
	\begin{align*}
		F(x) &= \int f(x) dx \\
		\pi &= \int_{-1}^{1} \frac{dx}{\sqrt{1-x^2}}
	\end{align*}

	\subsubsection{Equação com condições}

	\begin{equation*}
	f(x) = 
		\begin{cases}
			x^2 & \text{se } x \geq 0 \\\\
			-x & \text{se } x < 0
		\end{cases}
	\end{equation*}

	\subsection{Matrizes}
	
	\subsubsection{Parenteses}
	
	\begin{equation*}
		\begin{pmatrix}
		1 & 2 \\
		3 & 4
		\end{pmatrix}
	\end{equation*}
	
	\subsubsection{Colchetes}
	
	\begin{equation*}
		\begin{bmatrix}
			1 & 2 \\
			3 & 4 \\
			5 & 6
		\end{bmatrix}
	\end{equation*}
	
	\subsubsection{Barras}
	
	\begin{equation*}
		\begin{vmatrix}
			1 & 2 & 3 \\
			4 & 5 & 6
		\end{vmatrix}
	\end{equation*}
	
% ----- Fórmulas quimicas -----

	\section{Fórmulas Químicas}

	\subsection{Inline}

	Reação química: \ce{H2 + O2 -> H2O}

	\subsection{Em bloco}

	Molécula de cafeína:
	\begin{equation}
	\ce{C8H10N4O2}
	\end{equation}
	
% ----- Figuras -----
	\section{Figuras e posição}
	\begin{figure}[h!]
		\centering
		\includegraphics[width=0.4\textwidth]{example-image-a}
		\caption{Figura na posição 'h', diz para o \LaTeX{} colocar AQUI (here)}
		\label{fig:imagem}
	\end{figure}

	O parâmetro \verb|[width=0.4\textwidth]| faz a largura máxima ser 40\% da largura do texto.

	\begin{figure}[!t]
		\centering
		\includegraphics[clip, trim=2.8cm 0cm 2.8cm 0cm]{example-image-b}
		\caption{Figura na posição 'b', diz para o \LaTeX{} colocar abaixo na página (bottom)}
	\end{figure}

	O parâmetro \verb|[clip, trim=2.8cm 0cm 2.8cm 0cm]| corta a imagem 2.8cm da esquerda, 0cm de baixo, 2.8cm da direita e 0cm de cima.

	\begin{figure}[!b]
		\centering
		\includegraphics[height=0.4\textwidth, angle=-90]{example-image-c}
		\caption{Figura na posição 't', diz para o \LaTeX{}] colocar acima na página (top)}
	\end{figure}
	
	O parâmetro \verb|[angle=-90]| gira a imagem 90º no sentido horário.
	
	\begin{figure}[!p]
		\centering
		\includegraphics[angle=90, width=\textwidth]{example-image}
		\caption{Figura na posição 'p', diz ao \LaTeX{} colocar em uma página separada apenas de imagens}
	\end{figure}

	O parâmetro \verb|[angle=90, width=\textwidth]|, aliado com o posicionamento 'p', faz com que a imagem ocupe uma página inteira sozinha e esteja virada no sentido da pagina.

	\lipsum[1-2]

	\subsection{Mais de uma imagem por figura}
	
	\begin{figure}[h]
		\centering
		\begin{subfigure}[b]{0.45\textwidth}
			\includegraphics[width=\textwidth]{example-image-a}
			\caption{Figura a}
			\label{fig:1a}
		\end{subfigure}
		\hfill
		\begin{subfigure}[b]{0.45\textwidth}
			\includegraphics[width=\textwidth]{example-image-b}
			\caption{Figura b}
			\label{fig:1b}
		\end{subfigure}
		\caption{Comparação entre duas figuras}
		\label{fig:comparacao}
	\end{figure}

	As posições podem ser usadas para praticamente qualquer ambiente \LaTeX{}, sejam figura, tabelas, gráficos, etc.

	\lipsum[1-2]

% ----- Desenhos e diagramas -----

	\section{Desenhos e diagramas}

	\begin{center}
		\begin{tikzpicture}
			\begin{scope}[fill opacity=0.5]
				\fill[blue] (0,0) circle (2);
				\fill[red]  (2,0) circle (2);
			\end{scope}
			
			\draw (0,0) circle (2);
			\draw (2,0) circle (2);
			\node at (-0.0,1.5) {$A$};
			\node at (2.0,1.5)  {$B$};
			\node at (1,0)      {$A \cap B$};
		\end{tikzpicture}
	\end{center}

	Ou ainda:

	\begin{center}
		\begin{tikzpicture}[scale=2]
			
			\draw[thick] (-1,0) -- (1,0);
			\filldraw[black] (0,0) circle (0.02); 
			
			\draw[thick] (0,0) -- ({2*sin(30)}, {-2*cos(30)});
			
			\filldraw[gray!30] ({2*sin(30)}, {-2*cos(30)}) circle (0.15);
			
			\draw[dashed] (0,0) -- (0,-2.2);
			
			\draw[->] (0.5,0) arc (0:-30:0.5);
			\node at (0.3,0.15) {$\theta$};
			
			\draw[->, red, thick] ({2*sin(30)}, {-2*cos(30)}) -- ++(0,-1) node[right] {$\vec{P}$};
			\draw[->, blue, thick] ({2*sin(30)}, {-2*cos(30)}) -- (0,0) node[midway, below left] {$\vec{T}$};
			
			\node at ({2*sin(30)}, {-2*cos(30) - 0.3}) {m};
			
		\end{tikzpicture}
	\end{center}

% ----- Circuitos elétricos -----

	\section{Circuitos Elétricos}

	\begin{center}
		\begin{circuitikz}
			\draw (0,0) to[battery, l=12V] (0,2)
			to[R, l=1k$\Omega$] (2,2)
			to[bulb] (2,0)
			-- (0,0);
		\end{circuitikz}
	\end{center}

% ----- Tabelas -----
	\section{Tabelas}
	
	\begin{table}[h!]
		\caption{Exemplo completo de tabela com recursos do LaTeX}
		\label{tab:tabela}
		\renewcommand{\arraystretch}{1.4}
	
		\begin{tabularx}{\textwidth}{||l|c|r|X||}
			\hline
			\multicolumn{4}{|c|}{\textbf{Tabela de Demonstração}} \\
			\hline
			\textbf{Nome} & \textbf{Idade} & \textbf{Nota} & \textbf{Observações (coluna X)} \\
			\hline
			Ana & 21 & \cellcolor{ForestGreen} 9.5 & Excelente desempenho ao longo do semestre. Sempre participativa. \\
			\hline
			\multirow{2}{*}{João} & \multirow{2}{*}{22} & 7.0 & Teve dificuldades iniciais, mas apresentou evolução. \\
			& & 8.5 & Nota da recuperação. \\
			\hline
			Maria & 20 & \cellcolor{ForestGreen} 10 & Aluna destaque, liderou trabalhos em grupo. \\
			\hline
			Carlos & 23 & \cellcolor{RedOrange} 6.5 & Necessita melhorar participação e entrega de atividades. \\
			\hline
			\multicolumn{4}{|c|}{\textit{Legenda:} Nota máxima = 10.0} \\
			\hline
		\end{tabularx}
	\end{table}
	
	\lipsum[3]
	
% ----- Gráficos -----
	\section{Gráficos}
	
	\subsection{Gráficos de função}
	
	\begin{center}
		\begin{tikzpicture}
			\begin{axis}[
				axis lines=middle,
				xlabel={$x$},
				ylabel={$f(x)$},
				title={Gráfico de $f(x) = x^2$},
				grid=both
				]
				\addplot[domain=-3:3, samples=100, thick, blue] {x^2};
			\end{axis}
		\end{tikzpicture}
	\end{center}
	
	\subsection{Gráficos de dados}
	
	\begin{center}
		\begin{tikzpicture}
			\begin{axis}[
				xlabel={Tempo (s)},
				ylabel={Velocidade (m/s)},
				title={Dados experimentais},
				grid=major,
				only marks
				]
				\addplot[mark=*, red] coordinates {
					(0,0) (1,1.5) (2,3.8) (3,5.9) (4,7.1) (5,10)
				};
			\end{axis}
		\end{tikzpicture}
	\end{center}
	
	\subsubsection{Gráfico com dados em arquivo csv}
	
	\begin{center}
		\begin{tikzpicture}
			\begin{axis}[
				xlabel={Tempo (s)},
				ylabel={Velocidade (m/s)},
				title={Importando dados de um .csv},
				grid=major
				]
				\addplot table[
				col sep=comma,
				x=tempo,
				y=velocidade
				] {dados.csv};
			\end{axis}
		\end{tikzpicture}
	\end{center}
	
% ----- Blocos de código -----

	\section{Códigos}
	
	\subsection{Inline}
	
	Para códigos citados durante um texto, pode-se usar o \verb|p = input("está dia?")|
	
	\subsection{Em bloco}
	
	\begin{lstlisting}[language=Python]
p = input("esta dia?")
if p == sim:
	print("hallo world!")
else:
	print("vai dormir")
	\end{lstlisting}
	
% ----- Qr-codes -----
	
	\section{Qr-codes}
	
	\begin{center}
		\qrcode{https://youtu.be/dQw4w9WgXcQ?si=YOYCOUYIoobepFxJ}
	\end{center}
	
% ----- Bibliografia e Citações -----

	\section{Bibliografia e Citações}

	Citações a referências dentro do próprio documento podem ser feitas referenciando a imagem \ref{fig:comparacao}, a tabela \ref{tab:tabela} ou para a equação \eqref{eq:pi}.

	Para links externos pode-se usar um link como \href{https://youtu.be/dQw4w9WgXcQ?si=YOYCOUYIoobepFxJ}{Aqui}.
	
	Ou simplesmente a url como: \url{https://youtu.be/dQw4w9WgXcQ?si=YOYCOUYIoobepFxJ}, caso queira que a url apareça. (url's não obedecem a margem)
	
	Para citações bibliográficas pode-se usar esta, no fim da frase. \cite{latex}.
	
	Ou pode-se citar \textcite{einstein} no meio do texto.
	
	Ou ainda citar apenas o autor \citeauthor{einstein} ou o ano \citeyear{einstein}.
	
	\printbibliography
		
\end{document}











































